% \iffalse meta-comment
%
% Copyright (C) 2017- by Yanshuo Chu <yanshuoc@gmail.com>
%
% This file may be distributed and/or modified under the
% conditions of the LaTeX Project Public License, either version 1.3a
% of this license or (at your option) any later version.
% The latest version of this license is in:
%
% http://www.latex-project.org/lppl.txt
%
% and version 1.3a or later is part of all distributions of LaTeX
% version 2004/10/01 or later.
%
% \fi
%
% ^^A ======================================================================
% ^^A hithesis 的 szmidterm 类实现
% ^^A
% ^^A 深圳校区博士中期报告（hithesisartplus）的全部代码。本类只支持
% ^^A type=doctor、stage=midterm、campus=shenzhen 这一种组合。
% ^^A
% ^^A 各个类共用的文件头（\NeedsTeXFormat、\ProvidesExplClass 等）留在 hithesis.dtx，
% ^^A 共用的底层模块在 src/shared/，hithesis.ins 里每个 \file 先取头，再取实现。
% ^^A 子图题与固定词汇两个模块与 hit-report.dtx 共用，从那边取。
% ^^A ======================================================================
%
% ^^A ======================================================================
% ^^A Module szmidtermcls: hitsz-midterm-report 类主体（深圳博士中期报告专用）
% ^^A ======================================================================
% ^^A ======================================================================
% ^^A Module szmidterm-options: 文档类选项的声明与后处理（type/campus/fontset 等）（hitsz-midterm-report）
% ^^A ======================================================================
%    \begin{macrocode}
%<*szmidterm-options>
% \changes{v3.2a}{0000/00/00}{加通用选择器 \cs{hit@if@option@TF}，与 \file{hithesisbook.cls}
%   一致。}
% 读标志位的那族判断函数（\cs{hit@if@option@TF} 等）定义在
% \file{src/shared/hit-options.dtx}，各个类共用。这里只写本类特有的两个：
% \cs{hit@stage@name} 取「开题」或「中期」；按校区、学位、阶段组合出来的判断由
% \cs{hit@define@condition} 生成，一条一行，见下。参数从 expl3 之外传进来，
% 里面的空格照原样保留；替换旧写法时注意 \cs{fi} 后面被跳过的那个空格。
%
% \changes{v3.2a}{0000/00/00}{补上 \option{style/glue} 与 \cs{hit@glue@skip}、
%   \cs{hit@glue@font@size}，与 \file{hithesisbook.cls} 对齐。默认关，版面不变}
% \option{glue} 开时给标题前后距与行距一个伸缩量，关时是固定值。
%    \begin{macrocode}
\cs_new:Npn \hit@stage@name
  {
    \bool_if:NTF \g__hit_opening_bool
      { \hit@stage@opening } { \hit@if@option@T { midterm } { \hit@stage@midterm } }
  }
\hit@define@condition{weihai@or@shenzhen@master}{ ( \g__hit_weihai_bool || \g__hit_shenzhen_bool ) && \g__hit_master_bool }
\hit@define@condition{shenzhen@master@opening}{ \g__hit_shenzhen_bool && \g__hit_master_bool && \g__hit_opening_bool }
\hit@define@condition{harbin@master@opening@or@midterm}{ \g__hit_harbin_bool && \g__hit_master_bool
     && ( \g__hit_opening_bool || \g__hit_midterm_bool ) }
\hit@define@condition{harbin@master@opening}{ \g__hit_harbin_bool && \g__hit_opening_bool && \g__hit_master_bool }
%    \end{macrocode}
%
% 处理文档类选项
%    \begin{macrocode}

%% 处理文档类选项
%    \end{macrocode}
%
% 键名与变量名不同的字符串选项，\option{cjk-font} 这种带连字符的键要走这个。
%    \begin{macrocode}
\clist_map_inline:nn { raggedbottom , toc } { \hit@declare@boolean@option {#1} { false } }
\hit@declare@boolean@option{newtxmath}{false}
\hit@declare@boolean@option{glue}{false}

\hit@declare@choice@option{type}{}{bachelor,master,doctor}{}

\hit@declare@choice@option{campus}{}{shenzhen,weihai,harbin}{}
\bool_lazy_or:nnF
  { \bool_if_p:N \g__hit_weihai_bool }
  { \bool_if_p:N \g__hit_shenzhen_bool }
  { \bool_gset_true:N \g__hit_harbin_bool }

\hit@declare@choice@option{stage}{}{opening,midterm}{}
%    \end{macrocode}
% \changes{v3.2a}{0000/00/00}{设置参考文献列表中的作者是否全大写和最大显示数量}
% \changes{v3.2a}{0000/00/00}{设置参考文献列表标号的对齐方式}
%    \begin{macrocode}

\bool_new:N \g_bib_authorupper_bool
\bool_set_false:N \g_bib_authorupper_bool
\keys_define:nn { hit / class }
  {
    bibauthorupper .code:n =
      {
        \str_if_eq:nnT {#1} {full}
          { \bool_set_true:N \g_bib_authorupper_bool }
      } ,
  }

\int_new:N \g_bib_maxauthor_int
\int_set:Nn \g_bib_maxauthor_int {3}
\keys_define:nn { hit / class }
  {
    bibmaxauthor .code:n =
      {
        \int_set:Nn \g_bib_maxauthor_int {#1}
        \int_if_zero:nT { \g_bib_maxauthor_int }
          { \int_gdecr:N \g_bib_maxauthor_int }
      } ,
  }

\str_new:N \g_bib_labelalign_str
\keys_define:nn { hit / class }
  { biblabelalign .code:n = { \str_set:Nn \g_bib_labelalign_str {#1} } }

\hit@declare@string@option{fontset}
%    \end{macrocode}
%
% \changes{v3.2a}{0000/00/00}{西文字体独立成 \option{font} 选项}
% 西文字体，取值与 \file{hithesisbook.cls} 相同。
%    \begin{macrocode}
\hit@declare@string@option@as{font-en}{font@en}
%    \end{macrocode}
%
% \changes{v3.2a}{0000/00/00}{中文字体独立成 \option{cjk-font} 选项}
% 中文字体。留空时照旧把 \option{fontset} 转给 \pkg{ctex}，给了值就由模板自己设置。
%    \begin{macrocode}
\hit@declare@string@option@as{font-cn}{font@cn}
%    \end{macrocode}
%
% \changes{v3.2a}{0000/00/00}{\option{font} 与 \option{cjk-font} 改名
%   \option{font-en} 与 \option{font-cn}，与模板其余选项的中英后缀一致。旧名仍
%   生效，只在日志里提示}
%    \begin{macrocode}
\keys_define:nn { hit / class }
  {
    font     .tl_gset:c = { hit@font@en } ,
    cjk-font .tl_gset:c = { hit@font@cn } ,
  }
%    \end{macrocode}
%
% \changes{v3.2a}{0000/00/00}{新增 \option{math-font} 与 \option{math-style} 选项}
% 数学字体与数学排版风格，取值与 \file{hithesisbook.cls} 相同。
%    \begin{macrocode}
\hit@declare@string@option@as{math-font}{mathfont}
\hit@declare@string@option@as{math-style}{mathstyle}
\keys_define:nn { hit / class }
  {
    unknown .code:n = { \hit@pass@or@warn { ctexart } {#1} } ,
  }
\PassOptionsToPackage{no-math}{fontspec}

\ProcessKeyOptions [~hit~/~class~]

%    \end{macrocode}
%
% \changes{v3.2a}{0000/00/00}{artplus-options 的三段选项校验改用
%   \cs{bool\_lazy\_any:nF}、\cs{bool\_lazy\_or:nnF} 与 \cs{bool\_lazy\_all:nF}。
%   第三段原来写成三层嵌套、同一段报错文本重复三遍，现在合成一次判断。报错
%   文本按 CONTRIBUTING §4.6 留在 expl3 之外。}
% 本类只支持深圳校区博士的中期报告，三个条件缺一不可，原来靠三层嵌套分别报错，
% 三处文本相同，合成一次判断。
%    \begin{macrocode}
\cs_new_protected:Npn \hit@szmidterm@error@type
  {
    \ClassError { hithesisartplus }
      {
        \MessageBreak Please~specify~thesis~type~in~option:~
        \MessageBreak type=[bachelor~|~master~|~doctor]~
      }
  }
\cs_new_protected:Npn \hit@szmidterm@error@stage
  {
    \ClassError { hithesisartplus }
      {
        \MessageBreak Please~specify~stage~in~option:~
        \MessageBreak stage=<opening|midterm>~
      }
  }
\cs_new_protected:Npn \hit@szmidterm@error@combination
  {
    \ClassError { hithesisartplus }
      {
        \MessageBreak This~document~class~only~support~midterm~report~for~doctor
        in~shenzhen~campus
      }
  }
\bool_lazy_any:nF
  {
    { \bool_if_p:N \g__hit_bachelor_bool }
    { \bool_if_p:N \g__hit_master_bool }
    { \bool_if_p:N \g__hit_doctor_bool }
  }
  { \hit@szmidterm@error@type }
\bool_lazy_or:nnF
  { \bool_if_p:N \g__hit_opening_bool }
  { \bool_if_p:N \g__hit_midterm_bool }
  { \hit@szmidterm@error@stage }
\bool_lazy_all:nF
  {
    { \bool_if_p:N \g__hit_doctor_bool }
    { \bool_if_p:N \g__hit_midterm_bool }
    { \bool_if_p:N \g__hit_shenzhen_bool }
  }
  { \hit@szmidterm@error@combination }
%    \end{macrocode}
%
% \changes{v3.2a}{0000/00/00}{选项后处理（AutoFakeBold 与 fontset 透传）从 artpluscls 移入 artplus-options 模块}
%    \begin{macrocode}
\RequirePackage{ifthen}
\tl_if_empty:NTF \hit@font@cn
  {
    \PassOptionsToPackage { AutoFakeBold=2 } { xeCJK }
    \tl_if_empty:NF \hit@fontset
      { \PassOptionsToClass { fontset=\hit@fontset } { ctexart } }
  }
  { \PassOptionsToClass { fontset=none } { ctexart } }
%</szmidterm-options>
%<*szmidtermcls-load>
\LoadClass[a4paper,UTF8,zihao=-4,scheme=plain]{ctexart}
\bool_gset_true:N \g__hit_base_class_loaded_bool
%    \end{macrocode}
% \changes{v3.2a}{0000/00/00}{模块内容改由 \file{hithesis.ins} 按运行期顺序直接拼入 cls，不再生成 \file{modules/*.sty}}
%
% \changes{v3.2a}{0000/00/00}{LoadClass 之后的所有标准包加载、配置整体打包至 artplus-deps 模块；artpluscls 仅保留 LoadClass}
%    \begin{macrocode}
%</szmidtermcls-load>
%    \end{macrocode}
% ^^A ======================================================================
% ^^A Module szmidterm-deps: 第三方宏包集中加载与基础配置（含 hyperlink/geometry 内联）（hitsz-midterm-report）
% ^^A ======================================================================
%    \begin{macrocode}
%<*szmidterm-deps-a>
%    \end{macrocode}
%
% 各种宏包
%    \begin{macrocode}

%% 各种宏包
\RequirePackage{etoolbox}
%    \end{macrocode}
%
% \changes{v3.2a}{0000/00/00}{artplus-deps-a 的引擎检查改用 \cs{sys\_if\_engine\_xetex:F}，
%   与 art-deps-a、book-deps-a 保持一致}
%    \begin{macrocode}
\sys_if_engine_xetex:F
  { \ClassError { hithesis } { Please~use:~\MessageBreak~xelatex } { } }
\RequirePackage{xparse}

\hit@setup@math
%    \end{macrocode}
%
% \changes{v3.2a}{0000/00/00}{数学部分抽到 \file{src/shared/hit-math.dtx}}
% 数学宏包、定理机制与数学字体的设置见 \file{src/shared/hit-math.dtx}。
%
% 删除 \pkg{newtxtext}，由于最新一次更新 (\texttt{revision 62369}) 中与 \pkg{ctex} 宏包冲突，且无法解决，将正文字体替换为 TG TermesX 与 TG Heros.
% \changes{v3.0.18}{2022/05/02}{删除 \pkg{newtxtext}，由于最新一次更新 (\texttt{revision 62369}) 中与 \pkg{ctex} 宏包冲突，且无法解决，将正文字体替换为 TG Termes 与 TG Heros.}
% \changes{v3.0.19}{2022/05/05}{将 TG Termes 字体替换为 TG TermesX 字体，与 \pkg{newtxmath} 宏包更适配，删除 \option{Scale} 选项。}
% \changes{v3.2a}{0000/00/00}{英文字体三件套设置从 artpluscls 拆出至 artplus-fonts 模块}
%    \begin{macrocode}
%</szmidterm-deps-a>
%    \end{macrocode}
% ^^A ======================================================================
% ^^A Module szmidterm-fonts: 正文字号、中英文主字体与字号宏命令（\wuhao 等）（hitsz-midterm-report）
% ^^A ======================================================================
%    \begin{macrocode}
%<*szmidterm-fonts>
%    \end{macrocode}
%
% \changes{v3.2a}{0000/00/00}{字体层主体抽到 \file{src/shared/hit-fonts.dtx}，
%   这里只留调用与本类特有的部分}
%    \begin{macrocode}
\hit@setup@fonts
\cs_set:Npn \normalsize
  {
    \@setfontsize \normalsize { 12bp } { 15bp }
    \abovedisplayskip = 6pt
    \abovedisplayshortskip = 6pt
    \belowdisplayskip = \abovedisplayskip
    \belowdisplayshortskip = \abovedisplayshortskip
  }
%</szmidterm-fonts>
%<*szmidterm-deps-a>

\RequirePackage{graphicx}
\RequirePackage{pdfpages}
\includepdfset{fitpaper=true}
\RequirePackage{enumitem}       %使用enumitem宏包,改变列表项的格式
\RequirePackage{environ}
\hit@setup@footnote@packages
%    \end{macrocode}
%
% \pkg{xeCJK} 包不再自动加载 \pkg{xeCJKfntef}
% \changes{v3.0.17}{2022/02/23}{\pkg{xeCJK} 包不再自动加载 \pkg{xeCJKfntef}}
%    \begin{macrocode}
\RequirePackage{xeCJKfntef}
%    \end{macrocode}
%
% \pkg{newtx} 更新至 v1.7 版本后与 \pkg{ctex} 宏包冲突, 需要手动引入 \option{nofontspec} 选项. 判断方法为是否存在 v1.7 更新后的新宏包 \pkg{newtx.sty}
% \changes{v3.0.17}{2022/02/23}{\pkg{newtx} 更新至 v1.7 版本后与 \pkg{ctex} 宏包冲突, 需要手动引入 \option{nofontspec} 选项. 判断方法为是否存在 v1.7 更新后的新宏包 \pkg{newtx.sty}}
%    \begin{macrocode}

\RequirePackage{longtable}
\hit@setup@table
\hit@setup@footnote
\hit@setup@publist
\RequirePackage{booktabs}
%    \end{macrocode}
%
% 重定义 \pkg{booktabs} 宏包的默认线宽，绘制三线表时不再需要显式设置，以实现更好的内容与格式分离
% \changes{v3.0.20}{2022/05/06}{重定义 \pkg{booktabs} 宏包的默认线宽，绘制三线表时不再需要显式设置，以实现更好的内容与格式分离}
%    \begin{macrocode}
\heavyrulewidth=1.5pt
\lightrulewidth=1pt
\cmidrulewidth=0.5pt

%    \end{macrocode}
%
% \changes{v3.2a}{0000/00/00}{natbib/hyperref/url 配置从 artplus-deps 拆出至 artplus-hyperlink 模块（对齐 book-hyperlink/-art）}
%    \begin{macrocode}
%</szmidterm-deps-a>
%    \end{macrocode}
% ^^A ======================================================================
% ^^A Module szmidterm-hyperlink: hyperref/url 配置（PDF 书签、链接颜色）（hitsz-midterm-report）
% ^^A ======================================================================
%    \begin{macrocode}
%<*szmidterm-hyperlink>
\hit@setup@hyperlink
%</szmidterm-hyperlink>
%<*szmidterm-deps-b>

%    \end{macrocode}
%
% \changes{v3.2a}{0000/00/00}{geometry 配置从 artplus-deps 拆出至 artplus-geometry 模块（对齐 book-geometry/-art）}
%    \begin{macrocode}
%</szmidterm-deps-b>
%    \end{macrocode}
% ^^A ======================================================================
% ^^A Module szmidterm-geometry: 版芯尺寸与页边距（hitsz-midterm-report）
% ^^A ======================================================================
%    \begin{macrocode}
%<*szmidterm-geometry>
\RequirePackage{geometry}
\geometry{%根据PlutoThesis 原版定义而来
  a4paper, % 210 * 297mm
  hcentering,
  ignoreall,
  nomarginpar,
  centering,
  text={150true~mm,236true~mm},
  left=30true~mm,
  head=5true~mm,
  headsep=2true~mm,
  footskip=0true~mm,
  foot=5.2true~mm
}
%</szmidterm-geometry>
%<*szmidterm-deps-c>

\RequirePackage{tikzpagenodes}
\RequirePackage{eso-pic}
\RequirePackage{tikz}

\RequirePackage{fancyhdr}
\RequirePackage{tabularx}
\RequirePackage{xltabular}
\RequirePackage{varwidth}
\RequirePackage{changepage}
\RequirePackage{multicol}
\RequirePackage[below]{placeins}%允许上一个section的浮动图形出现在下一个section的开始部分,还提供\FloatBarrier命令,使所有未处理的浮动图形立即被处理
\RequirePackage{flafter}       % 使得所有浮动体不能被放置在其浮动环境之前，以免浮动体在引述它的文本之前出现.
\RequirePackage{multirow}       %使用Multirow宏包，使得表格可以合并多个row格
\RequirePackage{subcaption}%支持子图，必须在 ccaption 之前
\RequirePackage{ccaption} %支持双语标题
%    \end{macrocode}
%
% \changes{v3.2a}{0000/00/00}{\cs{bicaption} 支持省略英文图表名的四参数写法}
% \pkg{ccaption} 的 \cs{bicaption} 是五个参数，第四个是第二条题注用的图表名，
% 用户得写成 \texttt{Fig.} 加一个负细空格才能把 \cs{fnum@figure} 插的空格抵掉。
% 这里让第四个参数可以省略，省略时取 \option{const} 族里的 \option{figure-prefix-en}
% 与 \option{table-prefix-en}。五个参数的写法照旧——那才是 \pkg{ccaption} 的标准形式，
% 换到别的模板不用改。
%    \begin{macrocode}
\cs_new_eq:NN \__hit_bicaption:w \bicaption
\msg_new:nnn { hithesis } { bicaption-five-args }
  {
    \iow_char:N \\ bicaption~不再接受五个参数。 \\
    英文图表名请改用~\iow_char:N \\ hitsetup{const={figure-prefix-en=...}}~设置， \\
    调用处只写~\iow_char:N \\ bicaption[标签]{中文题}{英文题}。
  }
\cs_new_protected:Npn \__hit_bicaption:nnnnn #1#2#3#4#5
  {
    \tl_if_novalue:nTF {#1}
      { \__hit_bicaption:w {#2} {#3} {#4} {#5} }
      { \__hit_bicaption:w [#1] {#2} {#3} {#4} {#5} }
%    \end{macrocode}
%
% 旧的五参数写法会在这里剩下一个 \marg{英文题} 掉进正文。探一眼，撞见就报错，
% 免得静默排出来。题注后面正常不会紧跟花括号组。
%    \begin{macrocode}
    \peek_meaning:NT \c_group_begin_token
      { \msg_error:nn { hithesis } { bicaption-five-args } }
  }
\RenewDocumentCommand \bicaption { o +m +m +g }
  {
    \IfNoValueTF {#4}
      { \__hit_bicaption:nnnnn {#1} { } {#2} { \use:c { hit@ \@captype @prefix@en } } {#3} }
      { \__hit_bicaption:nnnnn {#1} {#2} {#3} { \use:c { hit@ \@captype @prefix@en } } {#4} }
  }
%    \end{macrocode}
% % \changes{v3.0.21}{2022/05/11}{删除 \pkg{xltxtra} 宏包, 来实现规范样式的上标。 \issue[125]}
%    \begin{macrocode}

\RequirePackage{lipsum}%random text
%    \end{macrocode}
%
% 添加字号、行距等设置
%    \begin{macrocode}
%</szmidterm-deps-c>
%<*szmidterm-deps-c>
%    \end{macrocode}
%
% 进行字号、行距等设置
% \changes{v3.2a}{0000/00/00}{章节标题格式从 artpluscls 拆出至 artplus-chapter 模块}
%    \begin{macrocode}
%</szmidterm-deps-c>
%    \end{macrocode}
% ^^A ======================================================================
% ^^A Module szmidterm-chapter: 章节标题格式（字号、缩进、间距、编号样式）（hitsz-midterm-report）
% ^^A ======================================================================
%    \begin{macrocode}
%<*szmidterm-chapter>
\RequirePackage{xparse}
\RequirePackage{tikz}

%% 进行字号、行距等设置
\ctexset{%
  section={
    afterindent=true,
    beforeskip={7.5bp},%上下空0.5行
    afterskip={7.5bp},
    format={\normalsize},
    aftername=\enspace,
    fixskip=true,
    break={},
  },
  subsection={
    afterindent=true,
    beforeskip={7bp},
    afterskip={7bp},
    format={\normalsize},
    aftername=\enspace,
    fixskip=true,
    break={},
  },
  subsubsection={
    afterindent=true,
    beforeskip={3bp},
    afterskip={3bp},
    format={\normalsize},
    aftername=\enspace,
    fixskip=true,
    break={},
  },
  paragraph/afterindent=true,
  subparagraph/afterindent=true
}

%    \end{macrocode}
%
% \changes{v3.2a}{0000/00/00}{\cs{section} 迁到 expl3（与 \file{hithesisbook.cls} 一致）}
%    \begin{macrocode}
\cs_set_eq:NN \hit@section \section
\RenewDocumentCommand \section { s o m s }
  {
    \IfBooleanF {#4}
      { \noindent \tikz [ overlay ] \draw ( -5mm , 0mm ) -- ( 155mm , 0mm ) ; }
    \IfBooleanTF {#1}
      {% \section*
        \IfNoValueTF {#2}
          { \hit@section * {#3} }
          { \hit@section * [#2] {#3} }
      }
      {% \section
        \IfNoValueTF {#2}
          { \hit@section {#3} }
          { \hit@section [#2] {#3} }
      }
  }
%</szmidterm-chapter>
%<*szmidterm-deps-c>
%    \end{macrocode}
%
% 中文封面
%    \begin{macrocode}
%</szmidterm-deps-c>
%    \end{macrocode}
% ^^A ======================================================================
% ^^A Module szmidterm-meta: 元信息字段（\hitsetup 接受的 key）（hitsz-midterm-report）
% ^^A ======================================================================
%    \begin{macrocode}
%<*szmidterm-meta>

%% 中文封面
%    \end{macrocode}
% \changes{v3.1d}{2025/03/03}{取消了深圳博士中期——即reportplus——中未使用到的参数cover-title-cn}
%    \begin{macrocode}
\hit@define@term{title-line-one-cn}
\hit@define@term{title-line-two-cn}
\hit@define@term{specialty-cn}
\hit@define@term{author-cn}
\hit@define@term{student-id}
\hit@define@term{class-id}
\hit@define@term{affiliation-cn}
\hit@define@term{supervisor-cn}
\hit@define@term{date-cn}
\hit@define@term{enroll-date-cn}

%    \end{macrocode}
%
% \changes{v3.2a}{0000/00/00}{关键词的逗号拆分与 \cs{hitsetup} 的入口挪到
%   \file{src/shared/hit-meta.dtx}}
% 关键词怎么拆、\cs{hitsetup} 的占位定义都见 \file{src/shared/hit-meta.dtx}。
%
% \changes{v3.2a}{0000/00/00}{\cs{hitsetup} 改用 l3keys 做入口，加一层按职责
%   划分的子族路由，与 \file{hithesisbook.cls} 一致。子族暂时只转发，键仍由
%   \texttt{hit/class} 族处理。}
% 入口换成 l3keys，根族路由到子族，子族一律把键转回
% \texttt{hit} 族。转发时值要重新加花括号，否则含逗号的值会被二次拆分成多个键。
%
% \changes{v3.2a}{0000/00/00}{元信息、可本地化字符串与参考文献选项正式注册进
%   对应的族，不再依赖子族兜底转发。}
% 元信息字段与可本地化字符串。做法与 \file{hithesisbook.cls} 相同：
% 前者的处理器直接调用同名命令，后者的键名由宏名去掉 \texttt{hit@} 前缀、
% 其余 \texttt{@} 换成连字符得来。
%    \begin{macrocode}
\hit@declare@info@field
  {
    affiliation-cn , author-cn , class-id , date-cn , enroll-date-cn , student-id , specialty-cn
    , supervisor-cn , cover-title-cn , title-line-one-cn , title-line-two-cn , report-type-cn
  }
\hit@declare@constant
  {
    authorization-usage-heading-cn , author-signature-field-label-cn , date-blank-cn ,
    date-field-label-front-cn , harbin-bachelor-school-bottom-mark , harbin-school-bottom-mark ,
    shenzhen-doctor-midterm-note , shenzhen-master-school-bottom-mark ,
    supervisor-signature-field-label-cn , bachelor-affiliation-field-label-cn ,
    bachelor-affiliation-field-label-shenzhen-cn , bachelor-affiliation-field-label-weihai-cn ,
    bachelor-class-field-label-cn , bachelor-date-field-label-cn ,
    bachelor-check-date-field-label-cn , bachelor-major-field-label-cn ,
    bachelor-student-id-field-label-cn , bachelor-student-field-label-cn ,
    bachelor-supervisor-field-label-cn , bachelor-document-type-cn , bachelor-degree-level-cn ,
    bachelor-teacher-comment-field-label-cn , bachelor-teacher-signature-field-label-cn ,
    affiliation-field-label-cn , associate-supervisor-field-label-cn , author-field-label-cn ,
    co-supervisor-field-label-cn , date-field-label-cn , degree-of-discipline-format-cn ,
    degree-field-label-cn , school-cn , school-field-label-cn , specialty-field-label-cn ,
    supervisor-field-label-cn , thesis-midterm-check-cn , document-type-cn ,
    thesis-title-prefix-cn , degree-level-cn , figure-prefix-en , table-prefix-en ,
    eqdenote-label-cn , eqdenote-label-en , eqdenote-dash , graduate-affiliation-field-label-cn
    , graduate-affiliation-major-field-label-cn , graduate-major-field-label-cn ,
    graduate-date-field-label-cn , graduate-enroll-date-field-label-cn ,
    graduate-student-id-field-label-cn , graduate-student-field-label-cn ,
    graduate-supervisor-field-label-cn , graduate-thesis-field-label-cn , hi ,
    classified-index-international-field-label-cn , classified-index-national-field-label-cn ,
    school-id , school-id-field-label-cn , secret-level-field-label-cn , shenzhen-campus ,
    stage-document-type , stage-midterm , stage-opening , title-separator-cn , weihai-campus
  }
%    \end{macrocode}
%
% 参考文献的三个选项，改用连字符命名。
%    \begin{macrocode}
\keys_define:nn { hit / bib }
  {
    label-align .code:n = { \str_set:Nn \g_bib_labelalign_str {#1} } ,
    author-case .code:n =
      { \str_if_eq:nnT {#1} { full } { \bool_set_true:N \g_bib_authorupper_bool } } ,
    max-author .code:n =
      {
        \int_set:Nn \g_bib_maxauthor_int {#1}
        \int_if_zero:nT { \g_bib_maxauthor_int } { \int_gdecr:N \g_bib_maxauthor_int }
      } ,
  }
\keys_define:nn { hit / style }
  {
    ragged-bottom .choice: ,
    ragged-bottom / true  .code:n = { \bool_gset_true:N \g__hit_raggedbottom_bool } ,
    ragged-bottom / false .code:n = { \bool_gset_false:N \g__hit_raggedbottom_bool } ,
    ragged-bottom .default:n = true ,
  }
%    \end{macrocode}
%
% \changes{v3.2a}{0000/00/00}{加旧名到新名的映射表与废弃警告。}
% 兼容期的处理，与 \file{hithesisbook.cls} 一致：旧名照常生效，只在日志里提示。
%    \begin{macrocode}
\prop_gset_from_keyval:Nn \g__hit_renamed_option_prop
  {
    raggedbottom   = { style = { ragged-bottom = ... } } ,
    biblabelalign  = { bib = { label-align = ... } } ,
    bibauthorupper = { bib = { author-case = ... } } ,
    bibmaxauthor   = { bib = { max-author = ... } } ,
  }
%    \end{macrocode}
%
% \changes{v3.2a}{0000/00/00}{元信息字段的旧名走废弃期：仍然可用，用一次提示一次}
% 旧名到新名的对照。旧名在 v3.1e 就是公开接口，写在用户的 \file{cover.tex} 里，
% 所以保留到 v4.1。
%    \begin{macrocode}
\hit@define@legacy@terms
  {
    ctitleone=title-line-one-cn ,
    ctitletwo=title-line-two-cn ,
    csubject=specialty-cn ,
    cauthor=author-cn ,
    cstudentid=student-id ,
    cclassid=class-id ,
    caffil=affiliation-cn ,
    csupervisor=supervisor-cn ,
    cdate=date-cn ,
    cenrolldate=enroll-date-cn
  }
\prop_gset_from_keyval:Nn \g__hit_renamed_load_option_prop
  {
    font     = { font-en = ... } ,
    cjk-font = { font-cn = ... } ,
  }
\__hit_check_class_options:
%    \end{macrocode}
%
% 根族路由。\option{toc} 既是族名又是老的布尔选项名，所以它单独处理：
% 值是 \texttt{true}/\texttt{false} 时当老写法，其余当子族的键值表。
%    \begin{macrocode}
\keys_define:nn { hit }
  {
    info  .meta:nn = { hit / info  } {#1} ,
    style .meta:nn = { hit / style } {#1} ,
    bib   .meta:nn = { hit / bib   } {#1} ,
    const .meta:nn = { hit / const } {#1} ,
    theorem .code:n = { \hit@theorem@style {#1} } ,
    toc .code:n =
      {
        \str_case:nnF {#1}
          {
            { true }  { \bool_gset_true:N \g__hit_toc_bool }
            { false } { \bool_gset_false:N \g__hit_toc_bool }
          }
          { \keys_set:nn { hit / toc } {#1} }
      } ,
    toc .default:n = true ,
    unknown .code:n = { \__hit_legacy_key:n {#1} } ,
  }
\RenewDocumentCommand \hitsetup { m } { \keys_set:nn { hit } {#1} }
%</szmidterm-meta>
%    \end{macrocode}
% ^^A ======================================================================
% ^^A Module szmidterm-pagestyle: 页眉页脚样式（hitsz-midterm-report）
% ^^A ======================================================================
% \changes{v3.2a}{0000/00/00}{页眉页脚样式从 artplus-frontmatter 的 \cs{makecover} 内联提取到独立的 artplus-pagestyle 模块（对齐 book-pagestyle/-art）}
%    \begin{macrocode}
%<*szmidterm-pagestyle>
\RequirePackage{fancyhdr}
\fancypagestyle{hit@artplus@main}{%
  \fancyhf{}%
  \cs_set:Npn \headrulewidth {0pt}%
  \cs_set:Npn \footrulewidth {0pt}%
  \fancyfoot[C]{\xiaowu-\nobreakspace{}\thepage\nobreakspace{}-}%
}
%</szmidterm-pagestyle>
%    \end{macrocode}
% ^^A ======================================================================
% ^^A Module szmidterm-frontmatter: 封面与扉页（中英、本科、博士后等多版本）（hitsz-midterm-report）
% ^^A ======================================================================
%    \begin{macrocode}
%<*szmidterm-frontmatter>

\cs_new_protected:Npn \hit@szmidterm@title@page@graduate
  {
~
 \xiaosi[1]\textcolor[rgb]{1,1,1}{\songti{\hit@hi}}%
 \vspace*{10mm}~
 \begin{center}~
 \songti\erhao\textbf{\hit@school@cn\hit@if@option@T{shenzhen}{\hit@shenzhen@campus}}
 \end{center}~
 \begin{center}~
 \songti\erhao\textbf{\hit@thesis@midterm@check@cn}
 \end{center}~
 \vspace{66bp}~
 \parbox[b][3cm][t]{\textwidth}{~
 \begin{center}\songti\sihao
 \cs_set:Npn \arraystretch {1.67}~
 \begin{tabular}{l@{\ \ }c}~
 \textbf{\hit@graduate@student@field@label@cn}~&~\textbf{\hit@author@cn}\\~
 \hline
 \textbf{\hit@graduate@enroll@date@field@label@cn}~&~\textbf{\hit@enroll@date@cn}\\~
 \hline
 \textbf{\hit@graduate@thesis@field@label@cn}~&~\textbf{\hit@title@line@one@cn}\\~
 \hline
 &\textbf{\hit@title@line@two@cn}\\~
 \hline
 \textbf{\hit@graduate@affiliation@major@field@label@cn}~&~\textbf{\hit@specialty@cn}\\~
 \hline
 \textbf{\hit@graduate@date@field@label@cn}~&~\textbf{\hit@date@cn}\\~
 \hline
 \textbf{\hit@graduate@supervisor@field@label@cn}~&~\textbf{\hit@supervisor@cn}\\~
 \hline
 \end{tabular}\cs_set:Npn \arraystretch {1}
 \end{center}~
 }~
  }

\cs_new_protected:Npn \hit@szmidterm@frame@picture
  {
%
~\begin{tikzpicture}[remember~picture,overlay]~
 \draw ([shift={(-5mm,0mm)}]current~page~text~area.south~west)~rectangle~([shift={(5mm,2mm)}]current~page~text~area.north~east);~
 \end{tikzpicture}%
  }
\cs_set_protected:Npn \makecover
  {
    \begin { titlepage }
      \hit@szmidterm@title@page@graduate
      \clearpage
      \hit@shenzhen@doctor@midterm@note
      \bool_if:NT \g__hit_toc_bool { \clearpage \tableofcontents \clearpage }
    \end { titlepage }
    \AddToShipoutPicture { \hit@szmidterm@frame@picture }
    \pagestyle { hit@artplus@main }
  }
%</szmidterm-frontmatter>
%<*szmidterm-deps-c>
%    \end{macrocode}
%
% 引用和参考文献格式
% \changes{v3.2a}{0000/00/00}{设置参考文献列表标号的对齐方式}
% \changes{v3.2a}{0000/00/00}{参考文献格式从 artpluscls 拆出至 artplus-bib 模块}
%    \begin{macrocode}
%</szmidterm-deps-c>
%    \end{macrocode}
% ^^A ======================================================================
% ^^A Module szmidterm-bib: 参考文献加载与 natbib/gbt7714 设置（hitsz-midterm-report）
% ^^A ======================================================================
%    \begin{macrocode}
%<*szmidterm-bib>
\RequirePackage[sort&compress,numbers]{natbib}
\RequirePackage{xparse}

%% 引用和参考文献格式
\cs_new_protected:Npn \bibstyle@numerical
  {
\bibpunct{[}{]}{,}{s}{,}{\textsuperscript{,}}
  }
\cs_new_protected:Npn \bibstyle@inline
  {
\bibpunct{[}{]}{,}{n}{,}{\hit@inline@separator}
  }

\citestyle{numerical}
\DeclareRobustCommand\inlinecite{\@inlinecite}
\cs_set:Npn \@inlinecite #1
  {
\group_begin: \citestyle{inline} \cs_set_eq:NN \@cite \NAT@citenum \citep{#1} \group_end:
  }
\cs_set_eq:NN \onlinecite \inlinecite

\cs_new_protected:Npn \hit@szmidterm@bibliography@heading
  {
    \section * { 参考文献 }
    \addcontentsline { toc } { section } { 参考文献 }
  }
\RenewDocumentEnvironment{thebibliography}{m}{%
  \list{\@biblabel{\@arabic\c@enumiv}}{%
    \str_if_eq:VnT \g_bib_labelalign_str {left} {
      \RenewDocumentCommand\makelabel{m}{##1\hfill}
    }
    \settowidth{\labelwidth}{\@biblabel{#1}}
    \dim_set:Nn \labelsep {0.5em}
    \dim_set:Nn \itemindent {0pt}
    \dim_set:Nn \leftmargin {\labelsep+\labelwidth}
    \dim_add:Nn \itemsep {-0.8em}
    \usecounter{enumiv}%
    \cs_set_eq:NN \p@enumiv \@empty
  \cs_set:Npn \theenumiv {\@arabic\c@enumiv}}%
  \sloppy\frenchspacing
  \flushbottom
  \clubpenalty0
  \@clubpenalty \clubpenalty
  \widowpenalty0%
  \interlinepenalty-50%
  \sfcode`\.\@m
}{}
%    \end{macrocode}
% \changes{v3.2a}{0000/00/00}{为上面的环境补充了一对空大括号}
% \changes{v3.2a}{0000/00/00}{控制参考文献列表中作者是否全大写和最大显示数量}
%    \begin{macrocode}
%% 用以控制参考文献列表中单个作者的格式
\exp_args:NNe \cs_new:Nn \g_bib_author_cs:n {
    \bool_if:NTF \g_bib_authorupper_bool {
        \exp_not:N \text_uppercase:n {#1}
    } {
        {#1}
    }
}
%% 用以输出参考文献列表中的作者
\NewDocumentCommand{\bibauthor}{m}{
  \int_zero:N \l_tmpa_int
  \exp_args:Ne \tl_map_inline:nn {
    \exp_args:Ne \tl_tail:n {#1}
  } {
    \int_compare:nNnT
      \l_tmpa_int = \g_bib_maxauthor_int {
      \prg_break:n {\tl_head:N {#1}}
    }
    \int_incr:N \l_tmpa_int
    \g_bib_author_cs:n {##1}
  }
  \prg_break_point:
}
%</szmidterm-bib>
%<*szmidtermcls-tail>
%    \end{macrocode}
% 杂项
% \changes{v3.2a}{0000/00/00}{misc-artplus 模块合并为 artpluscls 内联钩子}
%    \begin{macrocode}
\AtEndOfClass{\sloppy}
%</szmidtermcls-tail>
%    \end{macrocode}